\documentclass[11pt,oneside]{article}

\pdfminorversion=7
\usepackage{QuizStyle}

\hypersetup{
    pdftitle={20261001 Quiz for MATH102 Calculus II},
    pdfauthor={Jungwoo Kim},
    pdfsubject={Quiz for MATH102 Calculus II},
    pdfcreator={LaTeX}
}

\begin{document}

\quiztitle{20261001 Quiz}

\begin{problem}{1.5}
Find a linear function $\ell\colon\R^2\to\R$ satisfying
\[
    \ell(1,2)=3,
    \qquad
    \ell(2,3)=5.
\]
\end{problem}

\begin{solution}
By Theorem~1.2, every linear function $\ell\colon\R^2\to\R$ has the form
\[
    \ell(x,y)=px+qy
\]
for some $p,q\in\R$. The given values yield
\[
    p+2q=3,
    \qquad
    2p+3q=5.
\]
Thus $p=q=1$, and hence
\[
    \boxed{\ell(x,y)=x+y}.
\]
\end{solution}

\begin{problem}{1.6}
Let $\mathbf u=(u_1,u_2)$, $\mathbf v=(v_1,v_2)$, and $\mathbf w=(w_1,w_2)$ be vectors in $\R^2$, and let $a,b,c\in\R$. Using the definitions
\[
    \mathbf u+\mathbf v=(u_1+v_1,u_2+v_2),
    \qquad
    c\mathbf u=(cu_1,cu_2),
    \qquad
    -\mathbf u=(-u_1,-u_2),
\]
show the following properties.
\begin{problemsubparts}
    \item $\mathbf u+\mathbf v=\mathbf v+\mathbf u$.
    \item $\mathbf u+(\mathbf v+\mathbf w)=(\mathbf u+\mathbf v)+\mathbf w$.
    \item $c(\mathbf u+\mathbf v)=c\mathbf u+c\mathbf v$.
    \item $(a+b)\mathbf u=a\mathbf u+b\mathbf u$.
    \item $\mathbf u+(-\mathbf u)=\mathbf 0$.
\end{problemsubparts}
\end{problem}

\begin{solution}
Each identity follows componentwise from the definitions.
\begin{problemsubparts}
    \item
    \[
        \mathbf u+\mathbf v
        =(u_1+v_1,u_2+v_2)
        =(v_1+u_1,v_2+u_2)
        =\mathbf v+\mathbf u.
    \]
    \item
    \[
        \mathbf u+(\mathbf v+\mathbf w)
        =(u_1+v_1+w_1,u_2+v_2+w_2)
        =(\mathbf u+\mathbf v)+\mathbf w.
    \]
    \item
    \[
        c(\mathbf u+\mathbf v)
        =(cu_1+cv_1,cu_2+cv_2)
        =c\mathbf u+c\mathbf v.
    \]
    \item
    \[
        (a+b)\mathbf u
        =(au_1+bu_1,au_2+bu_2)
        =a\mathbf u+b\mathbf u.
    \]
    \item
    \[
        \mathbf u+(-\mathbf u)=(0,0)=\mathbf 0.
    \]
\end{problemsubparts}
\end{solution}

\begin{problem}{1.16}
Let
\[
    \mathbf u=(1,3),
    \qquad
    \mathbf v=(3,1).
\]
\begin{problemsubparts}
    \item Determine whether $\mathbf u$ and $\mathbf v$ are linearly independent.
    \item Express $(4,4)$ as a linear combination of $\mathbf u$ and $\mathbf v$.
    \item Express $(4,5)$ as a linear combination of $\mathbf u$ and $\mathbf v$.
\end{problemsubparts}
\end{problem}

\begin{solution}
\begin{problemsubparts}
    \item If $a\mathbf u+b\mathbf v=\mathbf 0$, then
    \[
        a+3b=0,
        \qquad
        3a+b=0.
    \]
    These equations give $a=b=0$. Thus $\mathbf u$ and $\mathbf v$ are linearly independent.

    \item Solving
    \[
        a+3b=4,
        \qquad
        3a+b=4
    \]
    gives $a=b=1$. Hence
    \[
        \boxed{(4,4)=\mathbf u+\mathbf v}.
    \]

    \item Solving
    \[
        a+3b=4,
        \qquad
        3a+b=5
    \]
    gives
    \[
        a=\frac{11}{8},
        \qquad
        b=\frac{7}{8}.
    \]
    Therefore
    \[
        \boxed{(4,5)=\frac{11}{8}\mathbf u+\frac{7}{8}\mathbf v}.
    \]
\end{problemsubparts}
\end{solution}

\begin{problem}{1.21}
Let $\mathbf u=(u_1,u_2)$, $\mathbf v=(v_1,v_2)$, and $\mathbf w=(w_1,w_2)$. Prove the following properties of the dot product.
\begin{problemsubparts}
    \item $\mathbf u\cdot(\mathbf v+\mathbf w)=\mathbf u\cdot\mathbf v+\mathbf u\cdot\mathbf w$.
    \item $\mathbf u\cdot\mathbf v=\mathbf v\cdot\mathbf u$.
\end{problemsubparts}
\end{problem}

\begin{solution}
Using the definition of the dot product,
\begin{problemsubparts}
    \item
    \begin{align*}
        \mathbf u\cdot(\mathbf v+\mathbf w)
        &=u_1(v_1+w_1)+u_2(v_2+w_2)\\
        &=(u_1v_1+u_2v_2)+(u_1w_1+u_2w_2)\\
        &=\mathbf u\cdot\mathbf v+\mathbf u\cdot\mathbf w.
    \end{align*}
    \item
    \[
        \mathbf u\cdot\mathbf v
        =u_1v_1+u_2v_2
        =v_1u_1+v_2u_2
        =\mathbf v\cdot\mathbf u.
    \]
\end{problemsubparts}
\end{solution}

\begin{problem}{1.25}
Let $\ell\colon\R^2\to\R$ be linear and satisfy
\[
    \ell(2,1)=3,
    \qquad
    \ell(1,1)=2.
\]
Find a vector $\mathbf c\in\R^2$ such that
\[
    \ell(\mathbf u)=\mathbf c\cdot\mathbf u
\]
for every $\mathbf u\in\R^2$.
\end{problem}

\begin{solution}
By Theorem~1.3, such a vector has the form $\mathbf c=(p,q)$. The given values yield
\[
    2p+q=3,
    \qquad
    p+q=2.
\]
Thus $p=q=1$, so
\[
    \boxed{\mathbf c=(1,1)}.
\]
\end{solution}

\begin{problem}{1.26}
Find the cosine of the angle between the vectors
\[
    \mathbf u=(1,2),
    \qquad
    \mathbf v=(3,1).
\]
\end{problem}

\begin{solution}
By Theorem~1.5,
\[
    \cos\theta
    =\frac{\mathbf u\cdot\mathbf v}{\|\mathbf u\|\,\|\mathbf v\|}.
\]
Here
\[
    \mathbf u\cdot\mathbf v=5,
    \qquad
    \|\mathbf u\|=\sqrt5,
    \qquad
    \|\mathbf v\|=\sqrt{10}.
\]
Therefore
\[
    \boxed{\cos\theta=\frac{5}{\sqrt{50}}=\frac{1}{\sqrt2}=\frac{\sqrt2}{2}}.
\]
\end{solution}

\begin{problem}{1.41}
Show that the vectors
\[
    (1,1,1,1),
    \quad
    (0,1,1,1),
    \quad
    (0,0,1,1),
    \quad
    (0,0,0,1)
\]
are linearly independent in $\R^4$.
\end{problem}

\begin{solution}
Let these vectors be $\mathbf v_1,\mathbf v_2,\mathbf v_3,\mathbf v_4$, respectively, and suppose
\[
    c_1\mathbf v_1+c_2\mathbf v_2+c_3\mathbf v_3+c_4\mathbf v_4=\mathbf 0.
\]
The first coordinate gives $c_1=0$. The second then gives $c_2=0$, the third gives $c_3=0$, and the fourth gives $c_4=0$. Hence the only linear combination equal to $\mathbf 0$ is the trivial one, so the four vectors are linearly independent.
\end{solution}

\begin{problem}{1.42}
Determine whether the vectors
\[
    (1,2,1),
    \quad
    (1,2,2),
    \quad
    (1,2,3),
    \quad
    (1,2,4)
\]
are linearly dependent or linearly independent. Identify the theorem of this section that is particularly applicable.
\end{problem}

\begin{solution}
There are four vectors in $\R^3$. By Theorem~1.8, any $n+1$ vectors in $\R^n$ are linearly dependent. Taking $n=3$, these four vectors are therefore
\[
    \boxed{\text{linearly dependent}}.
\]
The applicable result is Theorem~1.8.
\end{solution}

\begin{problem}{1.47}
Show that the vectors
\[
    (1,1,1),
    \qquad
    (1,2,3),
    \qquad
    (3,2,1)
\]
are linearly dependent.
\end{problem}

\begin{solution}
Let
\[
    \mathbf v_1=(1,1,1),
    \qquad
    \mathbf v_2=(1,2,3),
    \qquad
    \mathbf v_3=(3,2,1).
\]
We find the nontrivial relation
\[
    -4\mathbf v_1+\mathbf v_2+\mathbf v_3=\mathbf 0.
\]
Hence the vectors are linearly dependent.
\end{solution}

\begin{problem}{1.62}
Let
\[
    \mathbf c=(c_1,\dots,c_n)\in\R^n,
    \qquad
    \mathbf x=(x_1,\dots,x_n)\in\R^n.
\]
Show that the function
\[
    f(\mathbf x)=x_1+\mathbf c\cdot\mathbf x
\]
is linear by finding $\mathbf d=(d_1,\dots,d_n)$ such that
\[
    f(\mathbf x)=\mathbf d\cdot\mathbf x.
\]
\end{problem}

\begin{solution}
Since
\[
    f(\mathbf x)
    =(1+c_1)x_1+c_2x_2+\cdots+c_nx_n,
\]
take
\[
    \boxed{\mathbf d=(1+c_1,c_2,\dots,c_n)=\mathbf e_1+\mathbf c}.
\]
Then $f(\mathbf x)=\mathbf d\cdot\mathbf x$. The dot product with a fixed vector is linear in $\mathbf x$, so $f$ is linear.
\end{solution}

\begin{problem}{1.84}
Evaluate each of the following products.
\begin{problemsubparts}
    \item
    \[
        \begin{pmatrix}
            1&0&0\\
            0&2&0\\
            0&0&6
        \end{pmatrix}
        \begin{pmatrix}
            x_1\\x_2\\x_3
        \end{pmatrix}.
    \]
    \item
    \[
        \begin{pmatrix}
            3&-4\\
            4&3
        \end{pmatrix}
        \begin{pmatrix}
            6\\6
        \end{pmatrix}.
    \]
    \item
    \[
        \begin{pmatrix}
            \cos\theta&-\sin\theta\\
            \sin\theta&\cos\theta
        \end{pmatrix}
        \begin{pmatrix}
            1\\0
        \end{pmatrix}.
    \]
    \item
    \[
        \begin{pmatrix}
            \cos\theta&-\sin\theta\\
            \sin\theta&\cos\theta
        \end{pmatrix}^{2}.
    \]
\end{problemsubparts}
\end{problem}

\begin{solution}
Using the definition of matrix multiplication,
\begin{problemsubparts}
    \item
    \[
        \boxed{
        \begin{pmatrix}
            x_1\\2x_2\\6x_3
        \end{pmatrix}}.
    \]
    \item
    \[
        \begin{pmatrix}
            3&-4\\4&3
        \end{pmatrix}
        \begin{pmatrix}
            6\\6
        \end{pmatrix}
        =
        \boxed{
        \begin{pmatrix}
            -6\\42
        \end{pmatrix}}.
    \]
    \item
    \[
        \boxed{
        \begin{pmatrix}
            \cos\theta\\
            \sin\theta
        \end{pmatrix}}.
    \]
    \item
    \begin{align*}
        \begin{pmatrix}
            \cos\theta&-\sin\theta\\
            \sin\theta&\cos\theta
        \end{pmatrix}^{2}
        &=
        \begin{pmatrix}
            \cos^2\theta-\sin^2\theta&-2\sin\theta\cos\theta\\
            2\sin\theta\cos\theta&\cos^2\theta-\sin^2\theta
        \end{pmatrix}\\
        &=
        \boxed{
        \begin{pmatrix}
            \cos(2\theta)&-\sin(2\theta)\\
            \sin(2\theta)&\cos(2\theta)
        \end{pmatrix}}.
    \end{align*}
\end{problemsubparts}
\end{solution}

\begin{problem}{1.95}
Find an equation of each plane in the form
\[
    ax+by+cz=d.
\]
\begin{problemsubparts}
    \item The plane through the origin with normal vector $(1,0,0)$.
    \item The plane through $(0,0,0)$, $(0,1,1)$, and $(-3,0,0)$.
    \item The plane containing $(1,1,1)$ and parallel to the plane
    \[
        x-3y+5z=60.
    \]
\end{problemsubparts}
\end{problem}

\begin{solution}
\begin{problemsubparts}
    \item A normal vector is $(1,0,0)$, and the plane passes through the origin. Hence
    \[
        \boxed{x=0}.
    \]

    \item Two vectors in the plane are
    \[
        \mathbf v=(0,1,1),
        \qquad
        \mathbf w=(-3,0,0).
    \]
    Their cross product is
    \[
        \mathbf v\times\mathbf w=(0,-3,3),
    \]
    so a normal vector is $(0,-1,1)$. Since the plane passes through the origin,
    \[
        \boxed{-y+z=0}.
    \]

    \item A parallel plane has the same normal vector $(1,-3,5)$, so it has the form
    \[
        x-3y+5z=d.
    \]
    Substituting $(1,1,1)$ gives $d=3$. Therefore
    \[
        \boxed{x-3y+5z=3}.
    \]
\end{problemsubparts}
\end{solution}

\begin{problem}{1.104}
Using the formula
\[
    \mathbf v\times\mathbf w
    =(v_2w_3-v_3w_2,\,v_3w_1-v_1w_3,\,v_1w_2-v_2w_1),
\]
evaluate the following cross products.
\begin{problemsubparts}
    \item $(1,0,0)\times(0,1,0)$.
    \item $(0,1,0)\times(1,0,0)$.
    \item $(0,0,1)\times(a,b,c)$.
\end{problemsubparts}
\end{problem}

\begin{solution}
Substituting directly into the given formula,
\begin{problemsubparts}
    \item
    \[
        \boxed{(1,0,0)\times(0,1,0)=(0,0,1)}.
    \]
    \item
    \[
        \boxed{(0,1,0)\times(1,0,0)=(0,0,-1)}.
    \]
    \item
    \[
        \boxed{(0,0,1)\times(a,b,c)=(-b,a,0)}.
    \]
\end{problemsubparts}
\end{solution}

\begin{problem}{1.105}
Using
\[
    \mathbf e_1\times\mathbf e_2=\mathbf e_3,
    \qquad
    \mathbf e_2\times\mathbf e_3=\mathbf e_1,
    \qquad
    \mathbf e_3\times\mathbf e_1=\mathbf e_2,
\]
and the distributive and anticommutative laws
\[
    \mathbf u\times(\mathbf v+\mathbf w)
    =\mathbf u\times\mathbf v+\mathbf u\times\mathbf w,
    \qquad
    \mathbf v\times\mathbf w=-\mathbf w\times\mathbf v,
\]
evaluate the following cross products.
\begin{problemsubparts}
    \item $\mathbf e_1\times\mathbf e_2$.
    \item $\mathbf e_2\times(\mathbf e_1+\mathbf e_3)$.
    \item $(2\mathbf e_1+3\mathbf e_3)\times(a\mathbf e_1+b\mathbf e_2+c\mathbf e_3)$.
\end{problemsubparts}
\end{problem}

\begin{solution}
\begin{problemsubparts}
    \item
    \[
        \boxed{\mathbf e_1\times\mathbf e_2=\mathbf e_3}.
    \]
    \item
    \begin{align*}
        \mathbf e_2\times(\mathbf e_1+\mathbf e_3)
        &=\mathbf e_2\times\mathbf e_1+\mathbf e_2\times\mathbf e_3\\
        &=-\mathbf e_3+\mathbf e_1.
    \end{align*}
    Thus
    \[
        \boxed{\mathbf e_2\times(\mathbf e_1+\mathbf e_3)=\mathbf e_1-\mathbf e_3}.
    \]
    \item Using distributivity and anticommutativity,
    \begin{align*}
        &(2\mathbf e_1+3\mathbf e_3)\times(a\mathbf e_1+b\mathbf e_2+c\mathbf e_3)\\
        &\qquad=2b\mathbf e_3-2c\mathbf e_2+3a\mathbf e_2-3b\mathbf e_1\\
        &\qquad=-3b\mathbf e_1+(3a-2c)\mathbf e_2+2b\mathbf e_3.
    \end{align*}
    Hence
    \[
        \boxed{-3b\mathbf e_1+(3a-2c)\mathbf e_2+2b\mathbf e_3}.
    \]
\end{problemsubparts}
\end{solution}

\begin{problem}{2.4}
Show that a constant function $\mathbf F\colon\R^n\to\R^m$ is linear if and only if
\[
    \mathbf F(\mathbf x)=\mathbf 0
\]
for every $\mathbf x\in\R^n$.
\end{problem}

\begin{solution}
Suppose first that $\mathbf F$ is linear. Since $\mathbf F$ is constant,
\[
    \mathbf F(\mathbf x)=\mathbf F(\mathbf 0)
\]
for every $\mathbf x$. By linearity,
\[
    \mathbf F(\mathbf 0)
    =\mathbf F(0\mathbf x)
    =0\mathbf F(\mathbf x)
    =\mathbf 0.
\]
Thus $\mathbf F$ is the zero function.

Conversely, if $\mathbf F(\mathbf x)=\mathbf 0$ for every $\mathbf x$, then
\[
    \mathbf F(c\mathbf x)=\mathbf 0=c\mathbf F(\mathbf x)
\]
and
\[
    \mathbf F(\mathbf x+\mathbf y)=\mathbf 0
    =\mathbf F(\mathbf x)+\mathbf F(\mathbf y).
\]
Hence $\mathbf F$ is linear.
\end{solution}

\begin{problem}{2.5}
The plane in $\R^3$ given by
\[
    z=2x+3y
\]
is the $0$ level set of a linear function $\ell\colon\R^3\to\R$. Find such a function $\ell$.
\end{problem}

\begin{solution}
Rewrite the plane equation as
\[
    2x+3y-z=0.
\]
Thus the linear function
\[
    \boxed{\ell(x,y,z)=2x+3y-z}
\]
has exactly this plane as its $0$ level set.
\end{solution}

\begin{problem}{2.13}
Sketch or describe the $c$ level sets $f(\mathbf x)=c$ in $\R^3$ for the following functions $f\colon\R^3\to\R$ and values of $c$.
\begin{problemsubparts}
    \item $f(\mathbf x)=\|\mathbf x\|^2$, for $c=\frac12,1,2$.
    \item $f(\mathbf x)=\|\mathbf x\|$, for $c=\frac12,1,2$.
    \item $f(\mathbf x)=\dfrac{1}{\|\mathbf x\|^2}$, for $c=\frac12,1,2$.
    \item $f(\mathbf x)=\mathbf u\cdot\mathbf x$, for $c=\frac12,1,2$, where $\mathbf u$ is a unit vector. Hint: write $\mathbf x$ as the sum of a vector parallel to $\mathbf u$ and a vector orthogonal to $\mathbf u$.
\end{problemsubparts}
\end{problem}

\begin{solution}
\begin{problemsubparts}
    \item The equation $\|\mathbf x\|^2=c$ describes a sphere centered at the origin with radius $\sqrt c$. Thus the radii are
    \[
        \boxed{\frac{1}{\sqrt2},\ 1,\ \sqrt2}.
    \]

    \item The equation $\|\mathbf x\|=c$ describes a sphere centered at the origin with radius $c$. Thus the radii are
    \[
        \boxed{\frac12,\ 1,\ 2}.
    \]

    \item The equation
    \[
        \frac{1}{\|\mathbf x\|^2}=c
    \]
    is equivalent to $\|\mathbf x\|=1/\sqrt c$. Thus the level sets are spheres centered at the origin with radii
    \[
        \boxed{\sqrt2,\ 1,\ \frac{1}{\sqrt2}}.
    \]

    \item Write
    \[
        \mathbf x=t\mathbf u+\mathbf y,
        \qquad
        \mathbf y\perp\mathbf u.
    \]
    Since $\|\mathbf u\|=1$,
    \[
        \mathbf u\cdot\mathbf x=t.
    \]
    Hence the level set $\mathbf u\cdot\mathbf x=c$ is the plane perpendicular to $\mathbf u$ through the point $c\mathbf u$. For $c=\frac12,1,2$, the three planes pass through
    \[
        \boxed{\frac12\mathbf u,\ \mathbf u,\ 2\mathbf u},
    \]
    respectively.
\end{problemsubparts}
\end{solution}

\begin{problem}{2.35}
Show that the set
\[
    S=\{(x,y,z)\in\R^3\mid x^2+y^2<1\}
\]
is open.
\end{problem}

\begin{solution}
Let $\mathbf p=(a,b,c)\in S$ and set
\[
    \rho=\sqrt{a^2+b^2}<1,
    \qquad
    r=1-\rho>0.
\]
If $\mathbf q=(x,y,z)$ satisfies $\|\mathbf q-\mathbf p\|<r$, then by the triangle inequality in the $xy$-plane,
\begin{align*}
    \sqrt{x^2+y^2}
    &\leq \sqrt{(x-a)^2+(y-b)^2}+\sqrt{a^2+b^2}\\
    &\leq \|\mathbf q-\mathbf p\|+\rho\\
    &<r+\rho=1.
\end{align*}
Thus $\mathbf q\in S$, so $B_r(\mathbf p)\subseteq S$. Every point of $S$ is therefore an interior point, and $S$ is open by Definition~2.14.
\end{solution}

\begin{problem}{2.36}
Let
\[
    S=\R^2\setminus\{\mathbf 0\}.
\]
Show that $\mathbf 0$ is a boundary point of $S$.
\end{problem}

\begin{solution}
Let $r>0$. The open ball $B_r(\mathbf 0)$ contains $\mathbf 0$, which is not in $S$, and it also contains, for example,
\[
    \left(\frac r2,0\right)\in S.
\]
Thus every ball centered at $\mathbf 0$ contains a point of $S$ and a point not in $S$. By Definition~2.15, $\mathbf 0$ is a boundary point of $S$.
\end{solution}

\begin{problem}{2.37}
Let
\[
    T=\{(x,y)\in\R^2\mid x\geq0,\ y\geq0,\ x+y\leq1\}.
\]
\begin{problemsubparts}
    \item Describe the boundary of $T$.
    \item Show that the point $(0.0001,0.9998)$ is an interior point of $T$.
\end{problemsubparts}
\end{problem}

\begin{solution}
\begin{problemsubparts}
    \item The boundary consists of the three line segments
    \[
        \boxed{
        \partial T
        =\{(0,y)\mid 0\leq y\leq1\}
        \cup\{(x,0)\mid 0\leq x\leq1\}
        \cup\{(x,y)\mid x\geq0,\ y\geq0,\ x+y=1\}}.
    \]

    \item Let
    \[
        \mathbf p=(0.0001,0.9998),
        \qquad
        r=0.00005.
    \]
    If $\mathbf q=(x,y)$ satisfies $\|\mathbf q-\mathbf p\|<r$, then
    \[
        |x-0.0001|<r,
        \qquad
        |y-0.9998|<r.
    \]
    Therefore
    \[
        x>0.00005>0,
        \qquad
        y>0.99975>0,
    \]
    and
    \[
        x+y
        <0.0001+0.9998+2r
        =1.
    \]
    Hence $B_r(\mathbf p)\subseteq T$. By Definition~2.13, $\mathbf p$ is an interior point of $T$.
\end{problemsubparts}
\end{solution}

\end{document}
